\section{Discussion \& Observations}

\begin{enumerate}[leftmargin=1.2em,itemsep=2.0pt]
  \item \textbf{Elevated Tailing Factor, Injection~4:} Injection~4 showed a tailing factor of 1.97, compliant with the $\leq$2.0 acceptance limit but visibly elevated relative to the remaining five injections (range 1.28--1.31). Review of the pump pressure trace showed a transient 4\% pressure dip coincident with this injection, consistent with a small air bubble in the mobile phase line that resolved after re-priming. No corrective action was required, as the injection met the individual acceptance criterion, but the analyst flagged the observation in the instrument logbook per SOP-QC-014~\S7.3 for trend monitoring.
  \item \textbf{Resolution Margin:} The mean resolution between Impurity-A and the NXA-4471 main peak (3.43) provides a comfortable margin above the $\geq$2.0 requirement, indicating the method is robust to minor retention-time drift over the course of a sample analysis sequence.
  \item \textbf{Plate Count Trend:} Theoretical plate counts (6790--6910) are well above the 2000-plate minimum and consistent with historical qualification data for column Lot~C18-0552, indicating no loss of column efficiency since installation.
  \item \textbf{Impurity-B:} Impurity-B (RT~9.65\,min, RRT~$\approx$~1.23) was detected in the SST solution as expected from the reference material blend; it is not part of the critical resolution pair and carries no independent acceptance criterion under this method \cite{Dolan2002}.
\end{enumerate}
